\scanpage{94}
\mspage{93}

qui va impliquer, \deleted{\unsure{pour|que} \unsure{si}} avec Sous-at 3 (atelier plein) que si $\mathcal F$ satisfait At 9, $\mathcal F'$ aussi (\unsure{sans} que dans Sous at 4 il soit question des lieux).

Moyennant Sous at 4, \deleted{\unsure{At|4}} \inserted{et Sous at 3} At 10 pour $\mathcal F'$ implique \deleted{\unsure{que|pour}} \unsure{itou|item} pour $\mathcal F$.

At 11 bis \inserted{(\unsure{ou|en} At 11)} pour $\mathcal F'$ est une propriété de stabilité \inserted{de $\mathcal F'$ dans $\mathcal F$}, \unsure{nullement|nullmt} automatique. Idem pour At 12. Pour At 13, \deleted{\unsure{dans|de}} moyennant Sous at 3, Sous at 4, sa validité dans $\mathcal F$ l'implique dans $\mathcal F'$

At 14 est \unsure{abracadabra-|abracadabra|abracadabran}, At 14 bis dans la nature d'une condition de stabilité.

\begin{center}
\rule{1.35em}{0.4pt}
\end{center}
